\chapter{Estado da Arte}
\label{chap:estadodaarte}

\begin{introduction}
Este capítulo apresenta potenciais cenários (e a sua relevância) de monitorização do estado do utilizador durante o uso do computador, justificando a utilidade do uso dos periféricos do próprio computador ou sensores \textit{wearable off-the-shelf} em vez de sensores convencionais.
Adicionalmente, este capítulo analisa não só soluções de monitorização e integração de dados já existentes, mas também tecnologias relacionadas a \textit{streaming} que poderão ser utilizadas para este propósito.
\end{introduction}

\section{Monitorização do Estado do Utilizador durante o Uso do Computador}

\subsection{Monitorização da Atenção dos Alunos durante o Ensino Remoto}


Durante a pandemia de COVID-19, a necessidade de minimizar a progressão da doença levou à obrigatoriedade da realização de aulas de forma remota.

Esta transição, embora necessária, apresenta também problemas únicos que não existiam no estudo presencial, como a impossibilidade do docente monitorizar todos os alunos em simultâneo, especialmente durante a partilha do ecrã, o que se torna especialmente grave pois os alunos, tendo acesso ao telemóvel e ao próprio computador, têm uma maior tentação para se distrairem durante a aula~\cite{online_classes_sharma}.
Isto motiva a investigação de métodos de monitorização contínua dos alunos, complementares à supervisão do docente. Neste sentido, a integração de biosensores e dados multimodais (\textit{Multimodal Learning Analytics}) tem vindo a ganhar relevância como forma de captar estados cognitivos e emocionais dos alunos durante o ensino remoto~\cite{Becerra2026}. Contudo, de acordo com uma revisão sistemática realizada por Becerra \textit{et al.} (2026)~\cite{Becerra2026}, que analisou 54 estudos nesta área, a maioria da investigação existente ainda é realizada em ambientes laboratoriais controlados e com amostras pequenas e homogéneas — a mediana de participantes por estudo é de apenas 29{,}5 (média de 50{,}82), sendo que só um dos 54 estudos monitorizou mais de 100 alunos. Os próprios autores identificam esta limitação como uma barreira à validade externa e à generalização dos resultados, apontando a necessidade de estudos de maior escala e mais longitudinais, realizados em ambientes de aprendizagem reais.

Esta lacuna motiva o desenvolvimento de \textit{frameworks} abertas de monitorização capazes de monitorizar os alunos de forma não intrusiva durante o uso normal do computador, em vez de se limitarem a configurações experimentais controladas com um número reduzido de participantes, bem como a exploração do uso de fontes de dados de custo acessível ou de uso comum (como os próprios periféricos do computador).

\subsection{Monitorização em Ambiente Empresarial}

Com o crescente reconhecimento do impacto do stresse na produtividade dos trabalhadores e na sua saúde mental, existe um interesse crescente na monitorização contínua do stresse do trabalhador durante o dia de trabalho.
Historicamente, esta monitorização era realizada através de questionários subjetivos, contudo estes questionários não são fiáveis como método único de monitorização, por apenas medirem a perceção de stresse do próprio trabalhador (e não o seu estado real)~\cite{Kafkov2024}. 
Mais recentemente, com os avanços nas tecnologias \textit{wearable} e de Inteligência Artificial, o foco das investigações passou para a monitorização passiva utilizando \textit{wearables} ou dispositivos embarcados como os periféricos do próprio computador.

%TODO: falar do artigo de revisão


\section{Análise de Solução Existentes e Suas Limitações}

%TODO: apresentar as soluções existentes "similares" ao que se quer para os cenários acima, as suas características/funcionalidades/vantagens (se for descrito, a arquitetura/abordagem de recolha) e as limitações/desvantagens

%TODO: falar do gap que nenhuma destas preenche

\section{Tecnologias de Recolha e Processamento em \textit{Streaming}}

%TODO: definir brevemente o streaming como abordagem arquitetural/tecnológica (se calhar comparar com alternativas como batch)
%TODO: referir tecnologias de streaming/relacionadas que podem ser úteis, e as suas características (não é preciso dizer "esta ferramenta é X\% mais rápida que esta", focar-me no "porque é que as pessoas usam isto") 